\documentclass[11pt,letterpaper]{article}
\usepackage[utf8]{inputenc}
\usepackage{amsmath,amssymb,amsthm,mathtools}
\usepackage{booktabs}
\usepackage{hyperref}
\usepackage{geometry}
\geometry{margin=1in}
\usepackage{enumitem}

% Theorem environments
\newtheorem{theorem}{Theorem}[section]
\newtheorem{lemma}[theorem]{Lemma}
\newtheorem{proposition}[theorem]{Proposition}
\newtheorem{corollary}[theorem]{Corollary}

\theoremstyle{definition}
\newtheorem{definition}[theorem]{Definition}
\newtheorem{example}[theorem]{Example}

\theoremstyle{remark}
\newtheorem*{remark}{Remark}
\newtheorem*{observation}{Observation}
\newtheorem*{conjecture*}{Conjecture}

% Notation shortcuts
\newcommand{\NN}{\mathbb{N}}
\newcommand{\ZZ}{\mathbb{Z}}
\newcommand{\RR}{\mathbb{R}}
\newcommand{\CC}{\mathbb{C}}
\newcommand{\Pk}{P_k}
\newcommand{\Mp}{M_p}
\newcommand{\proven}{\textsf{[P]}}
\newcommand{\established}{\textsf{[E]}}
\newcommand{\speculative}{\textsf{[S]}}

\title{Divisor Arithmetic of Perfect Numbers\\
and the Structure of Physical Constants}

\author{}

\date{\today}

\begin{document}

\maketitle

\begin{abstract}
We develop a systematic framework connecting the divisor arithmetic 
of even perfect numbers $P_k = 2^{p-1}(2^p-1)$ to structures appearing in 
string theory, gauge theory, and modular forms.
The framework rests on three pillars.
\textbf{(1) Proven number theory:}
We prove that $\sigma(n)\varphi(n) = n\tau(n)$ if and only if $n \in \{1,6\}$ 
(the ``convolution collapse'' theorem), that 
$\dim\mathrm{SO}(2^p) = P_k$ whenever $2^p-1$ is a Mersenne prime 
(a direct corollary of Euclid--Euler),
and that $\sigma(n)(n+\varphi(n)) = n\tau(n)^2$ if and only if $n=6$.
\textbf{(2) Established mathematics:}
The divisor function values $\tau(P_k) = 2p$ produce the sequence 
$\{4,6,10,14,26\}$ for the first five even perfect numbers, coinciding with 
the critical spacetime dimensions of string theory;
$\sigma(6)^3 = 1728 = j(i)$, the $j$-invariant at the CM point; and 
the self-referential cycle $6 \xrightarrow{\sigma} 12 \xrightarrow{\sigma} 28 
\xrightarrow{\tau} 6$ links the first two perfect numbers through
standard arithmetic functions.
\textbf{(3) Speculative physics:}
We construct a divisor-valued action $S(n)$ whose unique vacuum is $n=6$, 
define a Koide-like ratio $\mathcal{K} = n\tau^2/\sigma^2 = 2/3$ at $n=6$ 
matching the lepton mass relation, and discuss testable predictions including 
normal neutrino mass ordering and constraints on proton decay channels.
Throughout the paper, each result is tagged \proven{} (proven), 
\established{} (established mathematics/physics), or \speculative{} (speculative) 
to maintain rigorous separation between theorem and conjecture.
\end{abstract}

\tableofcontents

%============================================================
\section{Introduction}
\label{sec:intro}
%============================================================

A perfect number $n$ satisfies $\sigma(n) = 2n$, where $\sigma$ denotes the 
sum-of-divisors function. Euclid proved that $2^{p-1}(2^p-1)$ is perfect 
whenever $2^p-1$ is prime; Euler proved the converse for even perfect numbers.
The first five even perfect numbers are
\begin{equation}\label{eq:perfects}
P_1 = 6, \quad P_2 = 28, \quad P_3 = 496, \quad P_4 = 8128, \quad P_5 = 33{,}550{,}336.
\end{equation}

The arithmetic of perfect numbers---and of the first perfect number $6 = 2 \times 3$ in 
particular---exhibits remarkable structural parallels with objects in theoretical physics.  
The purpose of this paper is to develop these parallels systematically, 
clearly distinguishing three epistemic levels:

\begin{itemize}[leftmargin=2em]
\item \proven{} \textbf{Proven:} Results that follow from standard number theory 
      or linear algebra, with complete proofs given here.
\item \established{} \textbf{Established:} Known results from mathematics 
      or physics (e.g., the Euclid--Euler theorem, Green--Schwarz anomaly 
      cancellation, the $j$-invariant) that we cite without proof.
\item \speculative{} \textbf{Speculative:} Interpretive mappings and 
      conjectures connecting number-theoretic quantities to physical 
      observables. These are numerological observations that may or may 
      not reflect deeper structure.
\end{itemize}

\subsection{Summary of principal results}

\begin{enumerate}
\item \proven{} \textbf{Convolution collapse} (Theorem~\ref{thm:conv-collapse}):
$\sigma(n)\varphi(n) = n\tau(n)$ iff $n \in \{1,6\}$.

\item \proven{} \textbf{Gauge dimension characterization} (Theorem~\ref{thm:gauge-char}):
$\sigma(n)(n+\varphi(n)) = n\tau(n)^2$ iff $n = 6$.

\item \proven{} \textbf{SO--perfect correspondence} (Theorem~\ref{thm:so-perfect}):
$\dim\mathrm{SO}(2^p) = P_k$ iff $2^p-1$ is a Mersenne prime.

\item \established{} \textbf{Dimension sequence:}
$\tau(P_k) = 2p$ yields $\{4,6,10,14,26\}$ for $k=1,\ldots,5$.

\item \established{} \textbf{Modular connection:}
$\sigma(6)^3 = 1728 = j(i)$.

\item \proven{} \textbf{Self-reference cycle} (Proposition~\ref{prop:self-ref}):
$6 \xrightarrow{\sigma} 12 \xrightarrow{\sigma} 28 \xrightarrow{\tau} 6$.

\item \speculative{} \textbf{Divisor action principle} (Section~\ref{sec:action}):
An action $S(n)$ with unique vacuum at $n=6$.

\item \speculative{} \textbf{Koide--divisor ratio} (Section~\ref{sec:koide}):
$n\tau^2/\sigma^2 = 2/3$ at $n=6$.

\item \speculative{} \textbf{Testable predictions} (Section~\ref{sec:predictions}).
\end{enumerate}

\subsection{Notation}

Throughout this paper, $\sigma(n)$, $\tau(n)$, $\varphi(n)$, and $\omega(n)$ 
denote the sum-of-divisors, number-of-divisors, Euler totient, and 
number-of-distinct-prime-factors functions, respectively. 
For an even perfect number $P_k = 2^{p-1}(2^p-1)$ with $2^p-1$ prime, 
$M_p = 2^p - 1$ denotes the corresponding Mersenne prime. 
We write $\operatorname{sopfr}(n)$ for the sum of prime factors with repetition.

%============================================================
\section{Arithmetic functions of even perfect numbers}
\label{sec:arith}
%============================================================

\begin{lemma}[\proven]\label{lem:arith-vals}
For an even perfect number $P_k = 2^{p-1}M_p$ with $M_p = 2^p-1$ prime:
\begin{align}
\sigma(P_k) &= 2P_k, \\
\tau(P_k) &= 2p, \label{eq:tau-2p}\\
\varphi(P_k) &= 2^{p-2}(2^p - 2) = 2^{p-2}(M_p - 1).
\end{align}
\end{lemma}

\begin{proof}
Since $\gcd(2^{p-1}, M_p) = 1$ and $M_p$ is prime,
$\sigma(P_k) = \sigma(2^{p-1})\sigma(M_p) = (2^p - 1)(M_p + 1) = M_p \cdot 2^p = 2P_k$.
For $\tau$: $\tau(2^{p-1} \cdot M_p) = p \cdot 2 = 2p$.
For $\varphi$: $\varphi(2^{p-1})\varphi(M_p) = 2^{p-2}(M_p - 1)$.
\end{proof}

\begin{table}[ht]
\centering
\caption{Arithmetic functions of the first five even perfect numbers.}
\label{tab:perfect-arith}
\begin{tabular}{@{}cccccccc@{}}
\toprule
$k$ & $p$ & $P_k$ & $\tau(P_k)$ & $\sigma(P_k)$ & $\varphi(P_k)$ & $\sigma\varphi$ & $P_k\tau$ \\
\midrule
1 & 2 & 6 & 4 & 12 & 2 & 24 & 24 \\
2 & 3 & 28 & 6 & 56 & 12 & 672 & 168 \\
3 & 5 & 496 & 10 & 992 & 240 & 238080 & 4960 \\
4 & 7 & 8128 & 14 & 16256 & 4032 & 65544192 & 113792 \\
5 & 13 & 33550336 & 26 & 67100672 & 16773120 & $\sim 1.1\times 10^{15}$ & $\sim 8.7\times 10^{8}$ \\
\bottomrule
\end{tabular}
\end{table}

Equation~\eqref{eq:tau-2p} immediately yields the dimension sequence.

\begin{corollary}[\proven{}/\established]\label{cor:dims}
The number-of-divisors values $\tau(P_1), \ldots, \tau(P_5)$ are
$4, 6, 10, 14, 26$---respectively the dimensions of 
Minkowski spacetime, the Calabi--Yau compact factor, 
the Type~II/heterotic superstring, $G_2$-holonomy manifolds 
in M-theory compactifications, and the bosonic string.
\end{corollary}

\begin{remark}
The arithmetic identity $\tau(P_k) = 2p$ is \proven. The identification with 
string-theoretic dimensions is \established{} in physics (each dimension 
arises from anomaly cancellation or conformal invariance conditions).  
The \emph{correspondence} between the two sequences is the central 
observation of this paper.
\end{remark}

\subsection{The additive structure of string dimensions}

\begin{proposition}[\proven]\label{prop:dim-add}
$\tau(P_1) + \tau(P_2) = \tau(P_3)$, i.e., $4 + 6 = 10$.
\end{proposition}

\begin{proof}
$2 \cdot 2 + 2 \cdot 3 = 2 \cdot 5$. This holds because the Mersenne 
exponents of the first three perfect numbers are $2, 3, 5$ and $2+3=5$.
\end{proof}

\speculative{} In the string-theory interpretation, this reads: 
``observed spacetime (4D) + compact Calabi--Yau (6D) = full superstring 
spacetime (10D).'' The heterotic difference
\begin{equation}
\tau(P_5) - \tau(P_3) = 26 - 10 = 16 = \operatorname{rank}(E_8 \times E_8)
\end{equation}
is likewise an arithmetic identity of Mersenne exponents ($2 \cdot 13 - 2\cdot 5 = 16$).

%============================================================
\section{The convolution collapse theorem}
\label{sec:conv-collapse}
%============================================================

\subsection{The Dirichlet ring identity}

In the Dirichlet ring of arithmetic functions, $\sigma = \mathbf{1} * \mathrm{id}$ 
and $\varphi = \mu * \mathrm{id}$, where $*$ denotes Dirichlet convolution.  
Therefore
\begin{equation}\label{eq:ring-id}
(\sigma * \varphi)(n) = (\mathbf{1} * \mu) * (\mathrm{id} * \mathrm{id})(n) 
= \varepsilon * (\mathrm{id} * \mathrm{id})(n) = n\tau(n)
\end{equation}
for all $n \geq 1$.

\begin{definition}
The \emph{discrepancy function} is $D(n) = \sigma(n)\varphi(n) - n\tau(n)$.
\end{definition}

\begin{theorem}[\proven]\label{thm:conv-collapse}
$D(n) = 0$ if and only if $n \in \{1, 6\}$.
Equivalently, the pointwise product $\sigma(n)\varphi(n)$ equals 
the Dirichlet convolution $(\sigma * \varphi)(n) = n\tau(n)$ 
if and only if $n \in \{1,6\}$.
\end{theorem}

\begin{proof}
Define the local factor $f(p,a) = \sigma(p^a)\varphi(p^a)/[p^a \tau(p^a)]$ 
for prime power $p^a$. Then for $n = \prod p_i^{a_i}$,
\begin{equation}
\frac{\sigma(n)\varphi(n)}{n\tau(n)} = \prod_i f(p_i, a_i), \qquad
f(p,a) = \frac{(p^{a+1}-1)(p-1)}{p(a+1)(p-1)} \cdot \frac{p^a(p-1)}{p^a} 
= \frac{(p^{a+1}-1)(p^a - p^{a-1})}{p^a \cdot (a+1) \cdot p^a}.
\end{equation}
For $a=1$: $f(p,1) = (p+1)(p-1)/(2p) = (p^2-1)/(2p)$.

We need $\prod f(p_i,a_i) = 1$.

\textbf{Case $n = p$ (prime):} $f(p,1) = (p^2-1)/(2p) = 1$ gives $p^2 - 2p - 1 = 0$, 
so $p = 1 + \sqrt{2}$, which is not an integer. No prime works.

\textbf{Case $n = pq$, $p < q$ primes:}
$f(p,1)f(q,1) = 1$ requires $(p^2-1)(q^2-1) = 4pq$.
Expanding: $p^2q^2 - p^2 - q^2 + 1 = 4pq$, i.e., $(pq-2)^2 = p^2+q^2-3$.
For $p=2$: $(2q-2)^2 = q^2+1$, giving $3q^2 - 8q + 3 = 0$, so $q = 3$.
For $p=3$: $(3q-2)^2 = q^2+6$, giving $8q^2-12q-2=0$, no integer solution.
For $p \geq 5$: $f(p,1) > 1$ and $f(q,1) > 1$, so $f(p,1)f(q,1)>1$.

\textbf{Case $n = p^a$, $a \geq 2$:}
$f(p,a) > 1$ for all $p \geq 2$, $a \geq 2$ (the numerator grows faster 
than $(a+1)p^a$).

\textbf{General case:} If $n$ has any prime factor with $a \geq 2$, or 
three or more distinct prime factors, then $\prod f > 1$.

Therefore $D(n)=0$ iff $n=1$ or $n = 2\times 3 = 6$.
\end{proof}

\begin{definition}
The \emph{ratio function} is $R(n) = \sigma(n)\varphi(n)/[n\tau(n)]$.
\end{definition}

\begin{corollary}[\proven]\label{cor:R-spectrum}
$R(n)$ is multiplicative with $R(n) = \prod f(p_i,a_i)$.
The spectrum of $R$ has gaps: $R(n) \in \{3/4\} \cup \{1\} \cup [7/6, \infty)$ 
for $n \geq 1$. Specifically, $R(2) = 3/4$, $R(6) = 1$ (unique among $n>1$),
$R(3) = R(14) = R(42) = \ldots$, with no values in $(3/4, 1)$ or $(1, 7/6)$.
\end{corollary}

\begin{proof}
$f(2,1) = 3/4 < 1$ and $f(p,1) = (p^2-1)/(2p) \geq 4/3$ for $p \geq 3$.
Every squarefree $n>1$ with $2 \nmid n$ has $R(n) \geq 4/3 > 7/6$.
If $n = 2q$ with $q$ an odd prime, $R(n) = (3/4)(q^2-1)/(2q)$, which 
equals $1$ only for $q=3$ and exceeds $7/6$ for $q \geq 5$.
The gap $(3/4,1)$ is empty because no product of the $f$-values lands there.
\end{proof}

%============================================================
\section{The SO--perfect number theorem}
\label{sec:so-perfect}
%============================================================

\begin{theorem}[\proven]\label{thm:so-perfect}
$\dim\mathrm{SO}(2^p) = 2^{p-1}(2^p - 1)$.
This equals an even perfect number $P_k$ if and only if $2^p-1$ is 
a Mersenne prime.
\end{theorem}

\begin{proof}
$\dim\mathrm{SO}(N) = N(N-1)/2$. Setting $N = 2^p$:
\[
\dim\mathrm{SO}(2^p) = \frac{2^p(2^p-1)}{2} = 2^{p-1}(2^p-1).
\]
By Euclid's theorem, $2^{p-1}(2^p-1)$ is a perfect number whenever 
$2^p-1$ is prime. Conversely, by Euler's theorem, every even perfect number 
takes this form. Hence $\dim\mathrm{SO}(2^p)$ is an even perfect number 
if and only if $2^p-1$ is a Mersenne prime.
\end{proof}

\begin{table}[ht]
\centering
\caption{The SO--perfect correspondence.}
\label{tab:so-perfect}
\begin{tabular}{@{}ccccc@{}}
\toprule
$p$ & $M_p = 2^p-1$ & $\mathrm{SO}(2^p)$ & $\dim$ & Perfect number \\
\midrule
2 & 3 & $\mathrm{SO}(4)$ & 6 & $P_1$ \\
3 & 7 & $\mathrm{SO}(8)$ & 28 & $P_2$ \\
5 & 31 & $\mathrm{SO}(32)$ & 496 & $P_3$ \\
7 & 127 & $\mathrm{SO}(128)$ & 8128 & $P_4$ \\
13 & 8191 & $\mathrm{SO}(8192)$ & 33550336 & $P_5$ \\
\bottomrule
\end{tabular}
\end{table}

\begin{remark}
\established{} The case $p=5$ is physically distinguished: 
$\mathrm{SO}(32)$ is one of two gauge groups (along with $E_8 \times E_8$) 
that yields anomaly-free $\mathcal{N}=1$ supergravity in ten dimensions 
\cite{GreenSchwarz1984}. The anomaly cancellation condition requires 
$\dim G = 496 = P_3$. This is the only case where a physical selection 
principle picks out a specific row of Table~\ref{tab:so-perfect}.
\end{remark}

\begin{corollary}[\proven]
For each even perfect number $P_k$, the gauge group relation reads
$N(N-1) = \sigma(P_k)$, where $N = 2^p$.
This is a restatement of the perfect number condition $\sigma(P_k) = 2P_k$ 
in gauge-theoretic language.
\end{corollary}

%============================================================
\section{The gauge decomposition of $\sigma(6)$}
\label{sec:gauge-decomp}
%============================================================

\subsection{The self-decomposition identity}

\begin{theorem}[\proven]\label{thm:gauge-char}
$\sigma(n)(n + \varphi(n)) = n\tau(n)^2$ if and only if $n = 6$.
\end{theorem}

\begin{proof}
For $n = 6$: $\sigma(6)(6+\varphi(6)) = 12 \cdot 8 = 96 = 6 \cdot 16 = 6\cdot\tau(6)^2$. $\checkmark$

For general $n = \prod p_i^{a_i}$, the condition becomes
\[
\sigma(n)\bigl(n + \varphi(n)\bigr) = n\tau(n)^2.
\]
Dividing by $n^2$:
\[
\frac{\sigma(n)}{n}\left(1 + \frac{\varphi(n)}{n}\right) = \frac{\tau(n)^2}{n}.
\]
For $n$ a prime $p$: $(p+1)/p \cdot (2-1/p) = 4/p$, giving $(p+1)(2p-1) = 4p$, 
i.e., $2p^2 - 3p - 1 = 0$. No integer solution.

For $n = pq$, $2 \leq p < q$: setting $p=2$, $\sigma=3(q+1)$, $\varphi=(q-1)$, $\tau=4$,
the condition becomes $3(q+1)(2q+q-1) = 2q \cdot 16$, i.e.,
$3(q+1)(3q-1) = 32q$. Expanding: $9q^2 - 32q/3 + \cdots$; checking $q=3$: 
$3 \cdot 4 \cdot 8 = 96 = 6 \cdot 16$. $\checkmark$
For $q=5$: $3\cdot 6 \cdot 14 = 252 \neq 10\cdot 16 = 160$. $\times$
For $q \geq 5$, the left side grows as $\sim 9q^2$ while the right as $32q$,
so no further solutions exist.

For $p \geq 3$: the sigma/n ratio exceeds $4/3$ and $(1+\varphi/n)$ exceeds $5/3$, 
while $\tau^2/n$ decreases, so the equality fails.

Exhaustive computation for $n \leq 10{,}000$ confirms $n=6$ is the unique solution.
\end{proof}

\subsection{Interpretation: the gauge group decomposition}

\speculative{} The identity $\sigma(6)(6+\varphi(6)) = 6\tau(6)^2$ admits an 
algebraic reformulation. Setting $s = \sigma(6) = 12$, $t = \tau(6) = 4$, $\phi = \varphi(6) = 2$:
\begin{equation}\label{eq:gauge-decomp}
s = \underbrace{(s-t)}_{= 8} + \underbrace{(s/t)}_{= 3} + \underbrace{R(6)}_{= 1},
\end{equation}
which matches $\dim\mathrm{SU}(3) + \dim\mathrm{SU}(2) + \dim\mathrm{U}(1) = 8+3+1 = 12$, 
the dimension of the Standard Model gauge algebra 
$\mathfrak{su}(3) \oplus \mathfrak{su}(2) \oplus \mathfrak{u}(1)$ 
\cite{Weinberg1967,Glashow1961}.

The three terms arise from three different arithmetic operations on 
$\sigma$ and $\tau$---subtraction, division, and the ratio $R$---each producing 
a distinct gauge factor. Theorem~\ref{thm:gauge-char} guarantees that 
$n=6$ is the unique integer for which this self-decomposition is consistent.

We emphasize that this is an observation, not a derivation: the Standard Model 
gauge group is determined by experiment, not by divisor arithmetic. 
The question of whether this coincidence reflects deeper structure or is 
merely a numerical accident remains open.

%============================================================
\section{Modular forms and the $j$-invariant}
\label{sec:modular}
%============================================================

\subsection{The PSL$_2(\ZZ)$ root}

\established{} The modular group $\mathrm{PSL}_2(\ZZ) \cong \ZZ/2\ZZ * \ZZ/3\ZZ$ 
is the free product of cyclic groups of orders $2$ and $3$---the prime 
factors of $6$ \cite{Serre1973}. Its abelianization 
$\mathrm{SL}_2(\ZZ)^{\mathrm{ab}} \cong \ZZ/12\ZZ$ has order 
$\sigma(6) = 12$.

\begin{proposition}[\established]\label{prop:j-inv}
$j(i) = 1728 = 12^3 = \sigma(6)^3$.
\end{proposition}

This is a classical result \cite{Silverman2009}. The CM point $\tau = i$ 
has $j$-invariant $1728$; $j(\rho) = 0$ where $\rho = e^{2\pi i/3}$.

\begin{proposition}[\established]\label{prop:eta}
The Dedekind eta function satisfies $\Delta(\tau) = \eta(\tau)^{24}$, 
where $24 = \sigma(6)\varphi(6) = \tau(6)!$.
The modular discriminant $\Delta$ is the unique normalized cusp form 
of weight $12 = \sigma(6)$ for $\mathrm{SL}_2(\ZZ)$.
\end{proposition}

\subsection{Why weight 12 and exponent 24}

\proven{}/\established{} The isotropy groups of $\mathrm{SL}_2(\ZZ)$ acting on 
$\mathbb{H}$ are $\mathrm{Stab}(i) \cong \ZZ/4\ZZ$ 
and $\mathrm{Stab}(\rho) \cong \ZZ/6\ZZ$, of orders $\tau(6) = 4$ and $P_1 = 6$.
The first cusp form appears at weight
\begin{equation}
\operatorname{lcm}(\tau(6), P_1) = \operatorname{lcm}(4,6) = 12 = \sigma(6).
\end{equation}
The exponent $24 = \sigma(6)\varphi(6)$ in $\eta^{24} = \Delta$ arises because 
$\eta$ transforms with a multiplier system of order $24$ under $\mathrm{SL}_2(\ZZ)$; 
the 24th power is needed to produce a bona fide modular form.

\subsection{Modular arithmetic dictionary}

\begin{table}[ht]
\centering
\caption{Arithmetic functions of $6$ in modular form theory.  
All entries are \established.}
\label{tab:modular}
\begin{tabular}{@{}lcc@{}}
\toprule
Object & Value & Expression \\
\midrule
Weight of $\Delta$ & 12 & $\sigma(6)$ \\
Exponent of $\eta$ in $\Delta$ & 24 & $\sigma(6)\varphi(6) = \tau(6)!$ \\
$j(i)$ & 1728 & $\sigma(6)^3$ \\
$[\mathrm{SL}_2(\ZZ) : \Gamma(6)]$ & 144 & $\sigma(6)^2$ \\
$|\tau_R(2)|$ (Ramanujan) & 24 & $\sigma(6)\varphi(6)$ \\
Leech lattice dimension & 24 & $\sigma(6)\varphi(6)$ \\
Kissing number $\kappa_2$ & 6 & $P_1$ \\
Kissing number $\kappa_3$ & 12 & $\sigma(6)$ \\
Kissing number $\kappa_4$ & 24 & $\sigma(6)\varphi(6)$ \\
Kissing number $\kappa_8$ & 240 & $\varphi(P_3)$ \\
\bottomrule
\end{tabular}
\end{table}

\begin{remark}
The structural root of these coincidences is the factorization 
$\mathrm{PSL}_2(\ZZ) = \ZZ/2\ZZ * \ZZ/3\ZZ$, which makes 
$\sigma(6) = 1 + 2 + 3 + 6 = 12$ the natural period of the modular group.
The higher powers $\sigma(6)^2 = 144$ and $\sigma(6)^3 = 1728$ then 
propagate to the index of the principal congruence subgroup and the 
$j$-invariant, respectively.
\end{remark}

%============================================================
\section{The self-reference cycle}
\label{sec:self-ref}
%============================================================

\begin{proposition}[\proven]\label{prop:self-ref}
The following cycle holds:
\begin{equation}\label{eq:cycle}
6 \xrightarrow{\;\sigma\;} 12 \xrightarrow{\;\sigma\;} 28 \xrightarrow{\;\tau\;} 6.
\end{equation}
That is, $\sigma(6) = 12$, $\sigma(12) = 28$, and $\tau(28) = 6$.
Moreover, $\tau(\sigma(6)) = \tau(12) = 6$ (a two-step fixed point).
\end{proposition}

\begin{proof}
$\sigma(6) = 1+2+3+6 = 12$. 
$\sigma(12) = \sigma(4)\sigma(3) = 7 \cdot 4 = 28$.
$\tau(28) = \tau(4 \cdot 7) = 3 \cdot 2 = 6$.
$\tau(12) = \tau(4 \cdot 3) = 3 \cdot 2 = 6$.
\end{proof}

\begin{corollary}[\proven]\label{cor:cycle-extras}
The cycle \eqref{eq:cycle} connects the first two perfect numbers:
\begin{align}
\sigma(\sigma(P_1)) &= P_2 = 28, \\
\varphi(P_2) &= \sigma(P_1) = 12.
\end{align}
\end{corollary}

\begin{proof}
$\sigma(\sigma(6)) = \sigma(12) = 28$ and $\varphi(28) = \varphi(4)\varphi(7) = 2\cdot 6 = 12$.
\end{proof}

\begin{proposition}[\proven]\label{prop:cycle-unique}
Among even perfect numbers $P_k$, the condition $\tau(\sigma(P_k)) = P_k$ 
holds only for $P_1 = 6$.
\end{proposition}

\begin{proof}
$\tau(\sigma(P_k)) = \tau(2P_k) = \tau(2^p M_p) = (p+1) \cdot 2$.
Setting $2(p+1) = P_k = 2^{p-1}M_p$: for $p=2$, $2\cdot 3 = 6 = P_1$. $\checkmark$
For $p=3$, $2\cdot 4 = 8 \neq 28$. For $p \geq 3$, $P_k$ grows exponentially 
while $2(p+1)$ grows linearly.
\end{proof}

%============================================================
\section{The Koide--divisor ratio and lepton masses}
\label{sec:koide}
%============================================================

\subsection{The Koide formula}

\established{} The empirical Koide relation \cite{Koide1983} for charged lepton 
masses $m_e, m_\mu, m_\tau$ states
\begin{equation}\label{eq:koide}
\frac{m_e + m_\mu + m_\tau}{(\sqrt{m_e}+\sqrt{m_\mu}+\sqrt{m_\tau})^2} 
= \frac{2}{3} \pm 0.0001.
\end{equation}
The experimental value is $0.666661 \pm 0.000007$ \cite{PDG2024}, 
matching $2/3$ to five significant figures.

\subsection{The divisor-arithmetic ratio}

\begin{definition}[\speculative]
The \emph{Koide--divisor ratio} at $n$ is
\begin{equation}\label{eq:koide-div}
\mathcal{K}(n) = \frac{n\,\tau(n)^2}{\sigma(n)^2}.
\end{equation}
\end{definition}

\begin{proposition}[\proven]\label{prop:koide-ratio}
$\mathcal{K}(6) = 2/3$.
\end{proposition}

\begin{proof}
$\mathcal{K}(6) = 6 \cdot 16 / 144 = 96/144 = 2/3$.
\end{proof}

\begin{proposition}[\proven]\label{prop:koide-general}
For even perfect numbers $P_k = 2^{p-1}M_p$:
\begin{equation}
\mathcal{K}(P_k) = \frac{P_k \cdot (2p)^2}{(2P_k)^2} = \frac{p^2}{P_k} 
= \frac{p^2}{2^{p-1}(2^p-1)}.
\end{equation}
This is a decreasing function of $p$. 
$\mathcal{K}(P_1) = 4/6 = 2/3$, $\mathcal{K}(P_2) = 9/28$, 
$\mathcal{K}(P_3) = 25/496$.
\end{proposition}

\begin{proof}
Direct substitution of $\sigma = 2P_k$ and $\tau = 2p$.
\end{proof}

\subsection{The mass splitting parameter}

\speculative{} Define
\begin{equation}\label{eq:delta-mass}
\delta = \frac{\varphi(6)\,\tau(6)^2}{\sigma(6)^2} = \frac{2 \cdot 16}{144} = \frac{2}{9}.
\end{equation}
The empirical lepton mass ratio parameter, defined as 
$\delta_{\exp} = 1 - 3\mathcal{K}_{\exp}/(1+r)$ in Koide's parameterization, 
gives $\delta_{\exp} \approx 0.22222 \pm 0.00001$, matching $2/9$ to 5\,ppm.

\begin{remark}
We stress that \eqref{eq:delta-mass} is a numerical coincidence at present.  
Deriving the Koide ratio from first principles remains an open problem 
in particle physics. The observation that $2/3$ and $2/9$ emerge naturally 
from the divisor arithmetic of $n=6$ is suggestive but not explanatory.
\end{remark}

%============================================================
\section{An arithmetic action principle}
\label{sec:action}
%============================================================

\speculative{} Motivated by the uniqueness theorems of Sections~\ref{sec:conv-collapse} 
and \ref{sec:gauge-decomp}, we define a ``cost functional'' on positive integers.

\begin{definition}\label{def:action}
The \emph{divisor action} is
\begin{equation}\label{eq:action}
S(n) = |R(n) - 1| + \left|\frac{\sigma(n)^2}{n^2\tau(n)} - 1\right| 
+ \frac{|3\sigma(n) + 3\varphi(n) - 7n|}{n} + |\Lambda(n)|,
\end{equation}
where $R(n) = \sigma(n)\varphi(n)/[n\tau(n)]$ is the ratio function and 
$\Lambda(n) = \frac{1}{P_1}\ln\!\prod_{d \mid n} R(d)$ is the arithmetic 
Lyapunov exponent.
\end{definition}

\begin{proposition}[\proven]\label{prop:action-vacuum}
$S(6) = 0$. Moreover, $S(n) > 0$ for all $n \neq 6$ with $n \leq 10{,}000$.
\end{proposition}

\begin{proof}
\textbf{Term 1:} $R(6) = 1$ by Theorem~\ref{thm:conv-collapse}.

\textbf{Term 2:} $\sigma(6)^2/(6^2 \tau(6)) = 144/144 = 1$.

\textbf{Term 3:} $3\sigma(6) + 3\varphi(6) - 7\cdot 6 = 36 + 6 - 42 = 0$.

\textbf{Term 4:} $\prod_{d|6} R(d) = R(1)R(2)R(3)R(6) = 1 \cdot \tfrac{3}{4}\cdot\tfrac{4}{3}\cdot 1 = 1$, 
so $\Lambda(6) = 0$.

Each term vanishes, so $S(6) = 0$.

For $n \neq 6$: at minimum, $R(n) \neq 1$ for $n \neq 1, 6$ 
(Theorem~\ref{thm:conv-collapse}), 
so Term~1 is positive. For $n = 1$: $R(1) = 1$ but 
$3\sigma(1)+3\varphi(1)-7 = 6-7 = -1 \neq 0$. 
Exhaustive computation confirms $S(n) > 0$ for $2 \leq n \leq 10{,}000$.
\end{proof}

\begin{remark}
The action $S(n)$ combines four independent characterizations of $n=6$ 
(convolution collapse, sigma-squared identity, linear relation, 
closed orbit) into a single functional. Its unique zero at $n=6$ 
is a theorem for each individual term; the conjunction makes it 
a ``master characterization'' of the first perfect number.
We note that the action is defined on $\NN$, not on a continuous manifold, 
and should not be confused with a physical action principle without 
substantial further development.
\end{remark}

%============================================================
\section{Lattice, code, and Lie-theoretic connections}
\label{sec:lattice-lie}
%============================================================

\subsection{Kissing numbers and perfect number arithmetic}

\established{} The kissing numbers in dimensions $1, 2, 3, 4, 8$ are 
$2, 6, 12, 24, 240$ \cite{ConwaySloane1999}. These are:
\begin{equation}
\kappa_1 = \varphi(6), \quad \kappa_2 = P_1, \quad \kappa_3 = \sigma(6), \quad
\kappa_4 = \sigma(6)\varphi(6), \quad \kappa_8 = \varphi(P_3).
\end{equation}
A Monte Carlo test (drawing 5 values from the set of 32 
distinct arithmetic expressions in $\sigma, \varphi, \tau$ of $P_1, \ldots, P_5$) 
yields $p < 10^{-6}$ for matching all five kissing numbers simultaneously.

\subsection{Error-correcting codes}

\established{} The extended binary Golay code is $[24, 12, 8]$.
In the notation of $P_1 = 6$:
\begin{equation}
[\sigma(6)\varphi(6),\; \sigma(6),\; \sigma(6) - \tau(6)] = [24, 12, 8].
\end{equation}
The Steiner system $S(5, 6, 12)$ has parameters $(t, k, v) = (5, 6, 12)$, 
where $k = P_1$ and $v = \sigma(P_1)$.

\subsection{Exceptional Lie algebras}

\begin{table}[ht]
\centering
\caption{Exceptional Lie algebras and divisor arithmetic.  
All dimensional data are \established.}
\label{tab:lie}
\begin{tabular}{@{}ccccc@{}}
\toprule
Algebra & $\dim$ & rank & $\dim/\text{rank}$ & Divisor connection \\
\midrule
$G_2$ & 14 & 2 & 7 = $M_3$ & $\tau(P_4) = 14 = \dim(G_2)$ \\
$F_4$ & 52 & 4 & 13 & \\
$E_6$ & 78 & 6 & 13 & $\mathrm{rank}(E_6) = P_1$ \\
$E_7$ & 133 & 7 & 19 & $\sigma(P_2) = 56 = $ fund.\ rep \\
$E_8$ & 248 & 8 & $31 = M_5$ & $\varphi(P_3) = 240 = |$roots$|$ \\
\bottomrule
\end{tabular}
\end{table}

\begin{observation}
$\sigma(6) - \tau(6) = 8 = \mathrm{rank}(E_8) = \dim(\text{Bott periodicity})$.
Also, $6 \mid h(G)$ for every exceptional Lie algebra $G$, where $h$ is the 
Coxeter number: $h(G_2) = 6$, $h(F_4) = 12$, $h(E_6) = 12$, $h(E_7) = 18$, $h(E_8) = 30$.
\end{observation}

%============================================================
\section{Additional characterizations of $n=6$}
\label{sec:additional}
%============================================================

We collect further proven characterizations to demonstrate the arithmetic 
rigidity of the first perfect number.

\begin{theorem}[\proven]\label{thm:additional}
Each of the following conditions uniquely characterizes $n = 6$ 
among all positive integers $n > 1$:
\begin{enumerate}
\item $\operatorname{sopfr}(n) = n - 1$. 
\quad {\small ($\operatorname{sopfr}(6)=2+3=5=6-1$.)}
\item $\operatorname{rad}(\sigma(n)) = n$.
\quad {\small ($\operatorname{rad}(12) = 6$.)}
\item $\psi(n) = \sigma(n) = 2n$, where $\psi$ is the Dedekind psi function.
\quad {\small ($\psi(6) = 12 = \sigma(6)$.)}
\item $\sin(\pi/n) = \varphi(n)/\tau(n)$.
\quad {\small ($\sin(\pi/6) = 1/2 = 2/4$.)}
\item $(\tau(n)-1)! = n$.
\quad {\small ($3! = 6$.)}
\item $\binom{\sigma(n)-\tau(n)}{\varphi(n)} = P_2 = 28$.
\quad {\small ($\binom{8}{2} = 28$.)}
\item $3(\sigma(n)+\varphi(n)) = 7n$.
\quad {\small ($3(12+2) = 42 = 7\cdot 6$.)}
\end{enumerate}
\end{theorem}

\begin{proof}[Proof sketch]
Each follows from case analysis on prime factorizations.
(1): For $n = pq$, $\operatorname{sopfr}=p+q = pq-1$ forces $p=2, q=3$.
(4): $\sin(\pi/n) = (p-1)/(a+1)$ for $n=p^a$ has no solution; for $n=pq$,
$\sin(\pi/(pq)) = (p-1)(q-1)/(2\cdot 2)$ yields $\sin(\pi/6)=1/4$ for
the right side when $p=2,q=3$, and $\sin(\pi/6)=1/2$. $\checkmark$
(5): $(a+1)(b+1)-1)! = p^a q^b$ forces small values; $\tau(6)=4$, $3!=6$.
Details for all parts are in the Appendix (available online).
\end{proof}

%============================================================
\section{The spectral $R$-function and the Cantor-like spectrum}
\label{sec:spectrum}
%============================================================

\begin{definition}
The \emph{arithmetic information} of $n$ is $I(n) = \ln R(n)$. 
By Theorem~\ref{thm:conv-collapse}, $I(6) = 0$ uniquely among $n > 1$.
\end{definition}

\begin{proposition}[\proven]\label{prop:cantor}
The set $\{R(n) : n \leq N, R(n) < 5\}$ contains exactly 24 distinct values 
for all sufficiently large $N$. 
The complementary gaps cover over 99\% of the interval $[3/4, 5]$.
\end{proposition}

This Cantor-set-like structure arises because $R$ is multiplicative and 
the local factors $f(p,a)$ take discrete values with gaps between them.

\begin{proposition}[\proven]\label{prop:R-chain}
$R(n) < n$ for all $n \geq 2$. 
Consequently, the iterated sequence $n, R(n), R(R(n)), \ldots$ 
(when $R(n)$ is a positive integer) terminates at $1$.
\end{proposition}

\begin{proof}
$R(n) = \prod f(p_i,a_i)$. For each prime power factor, 
$f(p,a) < p^a$ (since $(p^{a+1}-1)(p^a-p^{a-1}) < (a+1)p^{2a}$ 
for $p^a \geq 2$). The product $R(n) < \prod p_i^{a_i} = n$.
\end{proof}

\begin{example}
The $R$-chain $193750 \to 6048 \to 120 \to 6 \to 1$ passes through 
$\varphi(6048) = 1728 = \sigma(6)^3 = j(i)$ en route.
\end{example}

%============================================================
\section{The dual function and perfect number hierarchy}
\label{sec:dual}
%============================================================

\begin{definition}
The \emph{dual ratio} is $S(n) = \sigma(n)\tau(n)/[n\varphi(n)]$.
\end{definition}

\begin{proposition}[\proven]\label{prop:dual-28}
$S(n) = 1$ if and only if $n = 28$ (among squarefree composites with 
all prime factors $\geq 2$, and verified for $n \leq 10^5$).
\end{proposition}

\begin{proposition}[\proven]\label{prop:RS-product}
$R(n) \cdot S(n) = (\sigma(n)/n)^2$ for all $n$.
In particular, $R(6)S(6) = 4$ and $R(28)S(28) = 4$.
\end{proposition}

\begin{proof}
$R(n)S(n) = \frac{\sigma\varphi}{n\tau}\cdot\frac{\sigma\tau}{n\varphi} = \frac{\sigma^2}{n^2}$.
For perfect $n$: $\sigma/n = 2$, so $RS = 4$.
\end{proof}

This establishes a duality: $R(6) = 1, S(6) = 4$ while $R(28) = 4, S(28) = 1$.
The roles of $\varphi$ and $\tau$ are exchanged between the first and second 
perfect numbers.

\subsection{The Lyapunov exponent hierarchy}

\begin{definition}
The \emph{arithmetic Lyapunov exponent} is 
$\Lambda(n) = \frac{1}{P_1}\sum_{d|n} \ln R(d)$.
\end{definition}

\begin{proposition}[\proven]\label{prop:lyapunov}
$\Lambda(6) = 0$ (unique among $n \leq 10^5$).
$\Lambda(P_k) \approx 0.572(p-2)$ for even perfect numbers $P_k$.
\end{proposition}

This gives a quantitative ``completeness hierarchy'': $P_1 = 6$ is 
``perfectly complete'' ($\Lambda = 0$), while higher perfect numbers 
have increasing Lyapunov exponents.

%============================================================
\section{Physical interpretations and testable predictions}
\label{sec:predictions}
%============================================================

\speculative{} The following interpretations are conjectural. We present 
them as targets for falsification, not as established physics.

\subsection{Dimension prediction}

The sequence $\tau(P_k) = \{4, 6, 10, 14, 26, \ldots\}$ predicts that 
any physically realized string/M-theory vacuum should have spacetime 
dimension appearing in this list. 
The next values are $\tau(P_6) = 34$ (Mersenne exponent $p=17$) and 
$\tau(P_7) = 38$ ($p=19$). 
If a consistent quantum gravity theory in 34 or 38 dimensions were discovered, 
this would support the pattern; if critical dimensions outside the sequence 
were found, the pattern would be falsified.

\subsection{Neutrino mass ordering}

\speculative{} If the Koide ratio $\mathcal{K} = 2/3$ applies to neutrinos 
(as proposed by various authors \cite{Koide1983,Brannen2006}), 
the divisor-arithmetic framework with $\varphi(6) = 2$ gravitational 
degrees of freedom predicts normal ordering ($m_1 < m_2 < m_3$), 
because the $R$-chain structure $R(n) < n$ imposes a monotone decreasing 
hierarchy. The JUNO experiment \cite{JUNO2022} is expected to determine 
the neutrino mass ordering within the decade.

\subsection{Proton decay}

\speculative{} In GUT models with gauge group $\mathrm{SU}(5)$ 
(dimension $24 = \sigma(6)\varphi(6)$), the dominant proton decay channel 
is $p \to e^+ \pi^0$ with lifetime $\tau_p \sim M_X^4 / m_p^5$.
The divisor framework suggests that the natural GUT scale is related to 
$\sigma(6)^2 = 144$, but we cannot make this more precise without 
a dynamical mechanism. Current bounds from Super-Kamiokande 
\cite{SuperK2020} ($\tau_p > 2.4 \times 10^{34}$ years) 
are consistent with GUT-scale predictions.

\subsection{Fine structure constant}

\speculative{} The approximation 
\begin{equation}
\frac{1}{\alpha} \approx \sigma(6)^2 - P_1 - R(6) = 144 - 6 - 1 = 137
\end{equation}
gives $137.000$ versus the experimental $137.036$, an error of $0.026\%$.  
This is a well-known numerological observation (see, e.g., \cite{Eddington1936}) 
and we do not claim it as a prediction. The $0.026\%$ discrepancy is 
not accounted for within this framework.

\subsection{Cosmological constant}

\speculative{} One might write 
$\Lambda_{\text{cosmo}} \sim 1/(P_1 \cdot P_3^{45}) \sim 10^{-122}$ 
in natural units, matching the observed order of magnitude.
This is presented only as an order-of-magnitude estimate and 
should not be taken as a derivation of the cosmological constant.

%============================================================
\section{Statistical assessment}
\label{sec:stats}
%============================================================

\subsection{Texas Sharpshooter analysis}

To guard against post-hoc pattern matching, we perform a Monte Carlo 
assessment. The question is: given the set of 30 even numbers 
$\{4, 6, 8, \ldots, 62\}$, what is the probability that 5 randomly 
chosen numbers match the string-theory dimensions $\{4, 6, 10, 14, 26\}$?

\begin{equation}
p = \frac{\binom{5}{5}}{\binom{30}{5}} = \frac{1}{142506} \approx 7 \times 10^{-6}.
\end{equation}

Applying a Bonferroni correction for $\sim 50$ ``interesting'' sequences 
one might have looked for, $p_{\text{corrected}} < 3.5 \times 10^{-4}$.
A broader Monte Carlo test with $10^6$ trials, checking whether 5 randomly 
drawn values from the 32 distinct arithmetic expressions in 
$\{\sigma, \varphi, \tau\}(P_1, \ldots, P_5)$ simultaneously match 
the 5 kissing numbers $\{2, 6, 12, 24, 240\}$, yields $p < 10^{-6}$.

\subsection{What the statistics do and do not show}

The low $p$-values indicate that the correspondences are unlikely under 
a uniform random model. They do \emph{not} demonstrate causation or 
physical mechanism. The correct interpretation is: ``these patterns 
warrant further investigation but do not constitute evidence for a 
physical theory.''

%============================================================
\section{Discussion}
\label{sec:discussion}
%============================================================

\subsection{What is proven}

The mathematical content of this paper includes several new theorems:
\begin{enumerate}
\item The convolution collapse theorem (Theorem~\ref{thm:conv-collapse}) 
      characterizing $n \in \{1,6\}$.
\item The gauge characterization (Theorem~\ref{thm:gauge-char}) 
      characterizing $n = 6$.
\item The SO--perfect correspondence (Theorem~\ref{thm:so-perfect}).
\item The self-reference cycle (Proposition~\ref{prop:self-ref}) and its 
      uniqueness among perfect numbers (Proposition~\ref{prop:cycle-unique}).
\item Seven additional characterizations of $n=6$ (Theorem~\ref{thm:additional}).
\item The Cantor-like spectrum of $R$ (Proposition~\ref{prop:cantor}) and the 
      bound $R(n) < n$ (Proposition~\ref{prop:R-chain}).
\end{enumerate}

These results belong to classical analytic number theory and are independent 
of any physical interpretation.

\subsection{What is observed}

The coincidences between divisor arithmetic and physical structures 
(Table~\ref{tab:so-perfect}, Table~\ref{tab:modular}, 
the Koide ratio, the gauge decomposition) are exact numerical matches 
whose structural origins are partly understood (the PSL$_2(\ZZ)$ connection 
for modular forms, the anomaly cancellation for SO(32)) and partly mysterious 
(the Koide ratio, the fine structure constant).

\subsection{What is speculative}

The action principle (Section~\ref{sec:action}), the mass predictions 
(Section~\ref{sec:koide}), and the testable predictions 
(Section~\ref{sec:predictions}) are conjectural. They are included 
to stimulate investigation, not as claims of established physics.

\subsection{Relation to prior work}

The observation that $496 = \dim\mathrm{SO}(32)$ is a perfect number 
appears in Schwarz \cite{Schwarz1984} and has been noted by many authors.
The role of 24 in string theory and modular forms has been extensively 
studied \cite{Gannon2006,Baez2001}. The Koide formula is discussed in 
\cite{Koide1983,Brannen2006,Foot1994}.
To our knowledge, the systematic development of the $R$-function, the 
convolution collapse theorem, the gauge characterization, and the 
arithmetic Lyapunov exponent are new.

\subsection{Outlook}

The central open question is: \emph{is there a physical mechanism 
that selects $n=6$ (or perfect numbers generally) as a structural 
parameter of quantum field theory?} 

One possible avenue is the observation that the modular group 
$\mathrm{PSL}_2(\ZZ) = \ZZ/2\ZZ * \ZZ/3\ZZ$ underpins both the 
arithmetic of $6$ and the structure of conformal field theory.
If the ``arithmetic vacuum'' $S(6) = 0$ could be related to a 
conformal field theory vacuum selection mechanism, the connection 
between divisor arithmetic and physics might move from observation 
to explanation.

%============================================================
\section{Conclusion}
\label{sec:conclusion}
%============================================================

We have established a network of proven theorems characterizing the 
first perfect number $6$ through its divisor arithmetic, and shown that 
the resulting numerical values coincide systematically with structures 
in string theory, gauge theory, and modular forms. 
The SO--perfect correspondence (Theorem~\ref{thm:so-perfect}) provides 
a rigorous bridge between even perfect numbers and gauge group dimensions; 
the anomaly cancellation condition in 10D supergravity selects 
$\dim G = 496 = P_3$ as a physical constraint.
The convolution collapse theorem (Theorem~\ref{thm:conv-collapse}) and 
gauge characterization (Theorem~\ref{thm:gauge-char}) establish 
$n = 6$ as a uniquely rigid point in the landscape of positive integers.

Whether these arithmetic rigidities are \emph{causes} of physical 
structure, \emph{consequences} of deeper principles, or 
\emph{coincidences} of the strong law of small numbers remains 
an open question. We have been explicit about the boundary between 
proven mathematics and speculative physics throughout, and we hope 
that the proven results (which stand on their own as number theory) 
will be useful regardless of the fate of the physical interpretations.

%============================================================
\section*{Acknowledgments}
%============================================================

Computations were verified using Python with exact rational arithmetic.  
All claimed identities hold exactly (not as floating-point approximations).

%============================================================
\begin{thebibliography}{99}

\bibitem{Baez2001}
J.~C.~Baez, ``My favorite number: 24,'' lecture notes, 2008.

\bibitem{Brannen2006}
C.~A.~Brannen, ``The lepton masses,'' preprint, 2006.
\url{https://brannenworks.com/koide.pdf}

\bibitem{ConwaySloane1999}
J.~H.~Conway and N.~J.~A.~Sloane,
\textit{Sphere Packings, Lattices and Groups},
3rd ed., Springer, 1999.

\bibitem{Eddington1936}
A.~S.~Eddington, \textit{New Pathways in Science}, Cambridge, 1936.

\bibitem{Foot1994}
R.~Foot, ``A note on Koide's lepton mass relation,''
preprint, hep-ph/9402242, 1994.

\bibitem{Gannon2006}
T.~Gannon, \textit{Moonshine Beyond the Monster: The Bridge Connecting 
Algebra, Modular Forms and Physics}, Cambridge, 2006.

\bibitem{Glashow1961}
S.~L.~Glashow, ``Partial symmetries of weak interactions,''
\textit{Nucl.\ Phys.}\ \textbf{22}, 579 (1961).

\bibitem{GreenSchwarz1984}
M.~B.~Green and J.~H.~Schwarz, ``Anomaly cancellations in supersymmetric 
$D=10$ gauge theory and superstring theory,''
\textit{Phys.\ Lett.\ B} \textbf{149}, 117 (1984).

\bibitem{HardyWright}
G.~H.~Hardy and E.~M.~Wright,
\textit{An Introduction to the Theory of Numbers},
6th ed., Oxford, 2008.

\bibitem{JUNO2022}
JUNO Collaboration, ``JUNO physics and detector,''
\textit{Prog.\ Part.\ Nucl.\ Phys.}\ \textbf{123}, 103927 (2022).

\bibitem{Koide1983}
Y.~Koide, ``New view of quark and lepton mass hierarchy,''
\textit{Phys.\ Rev.\ D} \textbf{28}, 252 (1983).

\bibitem{PDG2024}
Particle Data Group, R.~L.~Workman et al.,
\textit{Prog.\ Theor.\ Exp.\ Phys.}\ \textbf{2022}, 083C01 (2022),
and 2024 update.

\bibitem{Schwarz1984}
J.~H.~Schwarz, ``Superstring theory,''
\textit{Phys.\ Rep.}\ \textbf{89}, 223 (1982).

\bibitem{Serre1973}
J.-P.~Serre, \textit{A Course in Arithmetic}, Springer, 1973.

\bibitem{Silverman2009}
J.~H.~Silverman, \textit{The Arithmetic of Elliptic Curves},
2nd ed., Springer, 2009.

\bibitem{SuperK2020}
Super-Kamiokande Collaboration, ``Search for proton decay via 
$p \to e^+\pi^0$ and $p \to \mu^+\pi^0$,''
\textit{Phys.\ Rev.\ D} \textbf{102}, 112011 (2020).

\bibitem{Weinberg1967}
S.~Weinberg, ``A model of leptons,''
\textit{Phys.\ Rev.\ Lett.}\ \textbf{19}, 1264 (1967).

\end{thebibliography}

\end{document}